\begin{agradecimentos}

    Agradeço, em primeiro lugar:

    Aos professores da Escola Municipal de Ensino Fundamental Bértolo Malacarne, do bairro de Cachoeira da Onça, no Município de São Garbiel da Palha, Espírito Santo.

    Aos professores do Instituto Federal de Educação, Ciência e Tecnologia do Espírito Santo, do campus de Nova Venécia.

    Aos professores da Faculdade de Ciência e Tecnologia em Engenharias da Universidade de Brasília,
    na cidade do Gama, Distrito Federal.

    Ao Professor Msc. Hilmer Rodrigues Neri, orientador deste trabalho.

    Pelo conhecimento, aprendizado e pelo suporte; não só a mim, mas a todos os que, como eu, tiveram a oportunidade de acesso à Educação. \emph{Viva à educação brasileira; pública, gratuita e de qualidade}.
   
    Agradeço, em segundo lugar:

    Aos meus pais, Alex Gomes de Oliveira e Geandra Storch, pelas palavras de afirmação, incentivo, carinho, orientação e suporte ao longo desta caminhada.

    Ao meu avô, Elias Storch, que me ensinou, com honestidade, paciência e lealdade, que o amor também canta baixinho na memória, como um bem-te-vi que anuncia sua presença.

    Aos meus amigos do Instituto Federal do Espírito Santo, após anos de amizade, companheirismo e apoio. Vocês fazem a diferença e sou extremamente grato por ter vocês em minha vida.

    Aos meus amigos da Universidade de Brasília, que me acompanharam na jornada do Ensino Superior e se tornaram a minha família no coração do Planalto Central.
    
    Agradeço, ainda, ao Distrito Federal por ter feito do cerrado brasileiro um local que eu sempre chamarei de lar. Hoje, tanto a minha história quanto a história do Brasil se dividem entre antes e depois de Brasília.

   Por fim, agradeço a todos que, de alguma forma, caminharam comigo até aqui. Levo comigo a gratidão por tudo o que vivi, por todos que encontrei e pela certeza de que esta jornada está apenas começando.

    Muito obrigado.
\end{agradecimentos}
